\section{Quels sont les apports de ce projet ?}

% ----------------------------------------------------------------
% 📌 À quoi sert cette section ?
% ----------------------------------------------------------------
% Cette section met en avant la valeur ajoutée du projet :
% - En quoi ce projet est-il bénéfique pour les étudiants porteurs ?
% - Quelles compétences ou connaissances seront développées ?
% - Quel est l’apport pour l’école, pour une entreprise, ou pour un utilisateur final ?
%
% Elle peut aborder des aspects techniques, pédagogiques, humains ou organisationnels.

% ----------------------------------------------------------------
% 🧭 Comment la remplir ?
% ----------------------------------------------------------------
% ➤ Listez les compétences que vous allez mobiliser ou acquérir (techniques et non techniques).
% ➤ Mentionnez les apports pour votre profil (développement, gestion, UX, communication...).
% ➤ Évoquez la valeur ajoutée pour SUPINFO, un partenaire, ou une communauté cible.
% ➤ Vous pouvez illustrer l’évolution ou l’enrichissement du profil de l’étudiant avec un tableau simple ou un schéma.

% ----------------------------------------------------------------
% 🧾 Exemple rédigé :
% ----------------------------------------------------------------
% Ce projet permettra aux étudiants de développer plusieurs compétences techniques, telles que :
% - La maîtrise de React pour la construction d’interfaces web interactives
% - L’utilisation de bibliothèques telles que DOMPurify pour sécuriser du contenu dynamique
% - La capacité à concevoir une architecture pédagogique structurée
%
% En parallèle, des compétences transversales seront renforcées :
% - Communication technique via la rédaction de documentation
% - Gestion de projet en mode agile (backlog, user stories)
% - Collaboration avec des enseignants ou partenaires externes
%
% Ce projet peut également devenir un outil pédagogique réutilisable à SUPINFO pour les promotions futures.

% ----------------------------------------------------------------
% 📊 Exemple de tableau d'apports (techniques / transversaux)
% ----------------------------------------------------------------

% \begin{table}[H]
% \centering
% \begin{tabular}{|l|p{8cm}|}
% \hline
% Domaine & Compétences développées \\
% \hline
% Technique & Programmation React, gestion des dépendances JS, sandbox sécurisé \\
% Pédagogique & Structuration de parcours d’apprentissage, gamification \\
% Transversal & Rédaction technique, gestion du temps, collaboration \\
% \hline
% \end{tabular}
% \caption{Synthèse des apports du projet}
% \end{table}

% ----------------------------------------------------------------
% ✍️ À vous de jouer : complétez ci-dessous
% ----------------------------------------------------------------

% (Début de réponse)
Ce projet permet aux porteurs de développer...

